\chapter{\nt{GLUT}与\nt{GABA}协同工作}
\label{sec:gaba}

\section{游戏规则}

仅由\nt{GLUT}神经元组成的网络不会选择、不能复位存储、也无法阻止活动发生雪崩，原因都在单调性（命题~\ref{prop:monotone}）。现在加入数量第二多的神经元类型——\nt{GABA}。规则不变：只用活体中真实存在的东西，不设时钟发生器。

神经元模型与~\eqref{eq:glutneuron} 相同，只是来自\nt{GABA}神经元的脉冲不是抬高、而是压低接收者的电压。权重现在有了两种符号，符号由发送者决定，即规则~\eqref{eq:dalesign}。

电路的几个参数。皮层中五分之一到四分之一的神经元是\nt{GABA}神经元~\cite{Hendry1987}。它们的输出很短：抑制的是邻居，而不是远处的脑区。抑制在一毫秒后到达，持续几毫秒。没有输入时\nt{GABA}神经元保持沉默：要抑制别人，它自己必须先受到兴奋，通常来自同一网络中的\nt{GLUT}神经元。

因此，从\nt{GLUT}神经元 $A$ 到神经元 $B$ 的负连接必须通过中介实现：$A$ 兴奋一个\nt{GABA}神经元，后者再抑制 $B$。改变符号的代价是一个神经元和一个延迟，约一毫秒。

对抑制而言，重要的不只是连接\emph{来自谁}，还有它\emph{落在哪里}：落点在很大程度上决定了连接的作用~\cite{Markram2004}。\nt{GLUT}输出主要落在接收者的输入分支即\emph{树突}上，携带数据：每个这样的输入是总和中的一项。落在神经元胞体上的输出作用于整个总和：许多快速\nt{GABA}神经元正是以这种方式对接收者的响应做除法，并决定它何时发放。落在接收者产生脉冲之处的输出是最后的裁决，一次否决整个输出。而落在别的导线末梢上的输出，只改变一条输入线路的释放概率 $p$，不影响其他线路~\cite{Rudomin1999}；第\ref{sec:pain}章中的疼痛闸门就是这样工作的。在脑外，输出落在肌纤维上（第\ref{sec:muscles}章）和器官上（第\ref{sec:chemoutput}章），还有一些输出哪里也不落，而是把递质释放到组织间隙或血液中（第\ref{sec:serocomputer}章和第\ref{sec:chemoutput}章）。

\section{非与否决}

现在能做非门了吗？差不多。输入沉默时还能工作的神经元，必须从某处获得兴奋。活体大脑中有这种兴奋：每个神经元不断收到来自成千上万其他神经元的背景活动。让背景使神经元 $Y$ 保持工作，输入 $x$ 经\nt{GABA}中介抑制它。这就是非门。

不过更有用的是\emph{否决}：“$a$，但 $b$ 出现时除外”：
\begin{empheq}[box=\keybox]{equation}
y \;=\; a \wedge \bar b .
\label{eq:veto}
\end{empheq}
神经元 $Y$ 从 $a$ 得到兴奋，从被 $b$ 兴奋的\nt{GABA}神经元得到抑制。这里不需要背景：兴奋由 $a$ 自己提供。用否决和或可以搭出一切。例如异或：$x\oplus y=(x\wedge\bar y)\vee(\bar x\wedge y)$，即两个否决神经元、两个\nt{GABA}中介和一个或神经元。

\begin{keyblock}
\begin{proposition}[\nt{GLUT}+\nt{GABA}的完备性]
\label{prop:gabacomplete}
由\nt{GLUT}和\nt{GABA}神经元组成的网络可以计算输入的任意逻辑函数，而且每个比特只用一根导线传送。第\ref{sec:glutcomputer}章中的“接通—断开”传感器不再需要。如果函数在所有输入都沉默时应输出1，就需要一个背景活动源。
\end{proposition}
\end{keyblock}

证明：只要有常数1，否决~\eqref{eq:veto} 加上或就是函数完备的，因为非 $x$ 就是否决“1，但 $x$ 出现时除外”。而在所有输入沉默时输出为0的函数，不需要常数1，只用否决和或就能搭出。

\section{运动方向}

时间上的否决，就像第\ref{sec:glutcomputer}章中时间上的与一样，给出运动方向检测器。

设两个相邻的传感器 $A$ 和 $B$ 相距 $\Delta x$。输出神经元 $Y$ 从 $B$ 得到兴奋，从 $A$ 经\nt{GABA}中介得到抑制，延迟为 $d$（图~\ref{fig:direction}）。以速度 $v$ 从 $A$ 移向 $B$ 的光斑在时刻 $0$ 经过 $A$，抑制在时刻 $d$ 到达 $Y$；光斑在时刻 $\Delta x/v$ 经过 $B$，兴奋也在这时到达。如果这两个时刻在窗口~\eqref{eq:window} 的精度内重合，抑制就抵消兴奋，$Y$ 保持沉默。这发生在速度约为
\begin{empheq}[box=\keybox]{equation}
v^* \;=\; \frac{\Delta x}{d}
\label{eq:vstar}
\end{empheq}
时。从 $B$ 移向 $A$ 时，来自 $B$ 的兴奋在时刻 $0$ 到达，而来自 $A$ 的抑制要到时刻 $\Delta x/v+d$ 才到，这时 $Y$ 已经发放了。

\begin{figure}[t]
\centering
\begin{tikzpicture}[>=Stealth,
  nrn/.style={circle,draw,very thick,fill=blue!6,minimum size=0.9cm,inner sep=0pt},
  gab/.style={nrn,fill=red!10},
  inh/.style={-{Bar[width=7pt]},thick,red!70!black}]
\node[nrn] (A) at (0,2) {$A$};
\node[nrn] (B) at (3,2) {$B$};
\node[gab] (G) at (0.9,0.6) {\small \nt{GABA}};
\node[nrn] (Y) at (3,-0.6) {$Y$};
\draw[->,thick] (A) -- (G);
\draw[inh] (G) -- node[below left=-2pt] {\scriptsize 延迟 $d$} (Y);
\draw[->,thick] (B) -- (Y);
\draw[->,very thick,red!70!black] (Y) -- (4.6,-0.6) node[right,black] {\small 输出};
\draw[->,thick,orange!85!black] (-0.6,2.9) -- (3.6,2.9) node[right,black] {\small 沉默};
\draw[<-,thick,orange!85!black] (-0.6,3.4) -- (3.6,3.4) node[right,black] {\small 响应};
\foreach \yy/\dd in {2.9/-1,3.4/1}{
  \foreach \k in {1,2,3}{\draw[orange!70!yellow,line width=0.8pt] (1.5+\dd*0.12+\dd*0.13*\k,\yy+0.1-0.1*\k+0.1) -- ++(\dd*0.25,0);}
  \fill[yellow!70!orange,draw=orange!70!black] (1.5,\yy) circle (0.14);}
\node[font=\scriptsize,align=center] at (1.5,3.95) {光斑};
\end{tikzpicture}
\caption{运动方向检测器。$B$ 兴奋 $Y$，$A$ 经\nt{GABA}中介以延迟 $d$ 抑制 $Y$。以约 $\Delta x/d$ 的速度从 $A$ 移向 $B$ 时，抑制抵消兴奋；反向运动时，抑制来迟了。带短线的黄色光斑表示在传感器上方移动的光。}
\label{fig:direction}
\end{figure}

在视网膜中，对运动方向的选择性看来正是这样产生的：来自不应引起响应的运动方向一侧的不对称抑制~\cite{Barlow1965}。

\section{减法还是除法}

到目前为止，我们的抑制做的是减法。实际上它的机制更精巧。

突触打开一个通道，也就是接入一个电导。兴奋性突触的平衡电位远高于阈值，比静息电位高约 $70\mV$：打开的通道把电压往上拉。抑制性\nt{GABA}突触的通道让氯离子通过，而氯的平衡电位在静息电位附近，所以打开的通道不是把电压往下拉，而是拉\emph{向静息}。以静息电位为零点，具有泄漏电导 $g_L$、兴奋性电导 $g_E$ 和抑制性电导 $g_I$ 的膜的稳态电压为
\begin{empheq}[box=\keybox]{equation}
V_\infty \;=\; \frac{g_E\,E_E}{g_L+g_E+g_I} ,
\label{eq:shunt}
\end{empheq}
其中 $E_E\approx70\mV$ 是兴奋性突触的平衡电位。这就是分压器的欧姆定律：$g_E$ 把电压拉向 $E_E$，$g_L$ 和 $g_I$ 把它拉向零。

$g_I$ 不是带着负号出现在分子里，而是出现在分母里：抑制不是从兴奋中\emph{减去}，而是把兴奋\emph{除以}某个数（图~\ref{fig:shunt}）。这种抑制称为分流抑制：打开的氯通道把膜短路了。对网络来说，这是增益调节。如果 $g_I$ 与邻居的总活动成正比，每个神经元响应的就不是自己输入的绝对强度，而是它在总活动中所占的份额。这种除以总活动的运算见于大脑的许多部位，被认为是大脑的标准运算之一~\cite{Carandini2012}：同一个网络在输入弱时和输入强时都能工作。

需要说明：公式~\eqref{eq:shunt} 讲的是不发放脉冲的神经元。对于正在发放的神经元，电压始终停在阈值附近，流过抑制性电导的电流几乎恒定，分流对脉冲\emph{频率}的作用主要是减法~\cite{HoltKoch1997}。只有同时给神经元加上相互平衡的兴奋性和抑制性背景电导，就像活体皮层那种抖动而平衡的状态，频率才会被除~\cite{Chance2002}。所以大脑的除法不是来自单独的分流，而是来自抑制加上背景噪声~\cite{Carandini2012}。

\begin{figure}[t]
\centering
\begin{tikzpicture}[>=Stealth,x=1.25cm,y=0.05cm]
\draw[->,gray!70!black] (0,0) -- (5.4,0) node[right,black] {\small $g_E/g_L$};
\draw[->,gray!70!black] (0,0) -- (0,68) node[above,black] {\small $V_\infty$，mV};
\foreach \v in {20,40,60}{\draw[gray!60!black] (-0.05,\v) -- (0.05,\v) node[left=3pt,font=\scriptsize] {\v};}
\foreach \x in {1,2,3,4,5}{\draw[gray!60!black] (\x,-1.5) -- (\x,1.5) node[below=5pt,font=\scriptsize] {\x};}
\draw[very thick,blue!60!black] plot[domain=0:5,samples=60] (\x,{70*\x/(1+\x)});
\draw[very thick,red!70!black] plot[domain=0:5,samples=60] (\x,{70*\x/(3+\x)});
\draw[thick,densely dashed,gray!50!black] plot[domain=0.56:5,samples=60] (\x,{70*\x/(1+\x)-25});
\node[right,font=\small,blue!60!black] at (5.02,58.3) {无抑制};
\node[right,font=\small,red!70!black] at (5.02,45.5) {除法：$g_I=2g_L$};
\node[right,font=\small,gray!40!black] at (5.02,32.5) {减去 $25\mV$};
\end{tikzpicture}
\caption{分流抑制对电压做除法，而不是减法。蓝色曲线是无抑制时的稳态电压~\eqref{eq:shunt}，红色曲线是加上抑制性电导 $g_I=2g_L$ 后的情形。红色曲线就是蓝色曲线沿水平方向拉伸三倍：仿佛输入被除以 $1+g_I/g_L=3$。虚线是减去恒定 $25\mV$ 的抑制：整条曲线下移，斜率不变。}
\label{fig:shunt}
\end{figure}

还有第二个推论。膜时间常数 $\tau_{\text{eff}}=C/(g_L+g_E+g_I)$，抑制会使它变短。而符合窗口~\eqref{eq:window} 与 $\tau$ 成正比，所以与兴奋同时到达的抑制会\emph{收窄窗口}。输入直接兴奋神经元，同时经\nt{GABA}中介以一毫秒的延迟抑制它，这是大脑中最常见的电路之一：神经元只来得及响应最初一两毫秒内到达的输入~\cite{Pouille2001}。

\section{选择}

现在来看\nt{GLUT}网络最缺的东西：从几个选项中选一个。取两群\nt{GLUT}神经元，输入分别为 $I_1$ 和 $I_2$。每群经自己的\nt{GABA}中介以强度 $\beta$ 抑制另一群（图~\ref{fig:wta}）。对频率 $r_1,r_2$ 有
\begin{equation}
\tau\,\frac{\dd r_1}{\dd t} = -r_1 + \bigl[I_1-\beta r_2\bigr]_+ ,\qquad
\tau\,\frac{\dd r_2}{\dd t} = -r_2 + \bigl[I_2-\beta r_1\bigr]_+ ,
\label{eq:wta}
\end{equation}
其中 $[u]_+=\max(u,0)$：频率不能为负。

\begin{figure}[t]
\centering
\begin{tikzpicture}[>=Stealth,
  nrn/.style={circle,draw,very thick,fill=blue!6,minimum size=0.9cm,inner sep=0pt},
  gab/.style={nrn,fill=red!10,minimum size=0.75cm},
  inh/.style={-{Bar[width=7pt]},thick,red!70!black}]
\node[nrn] (E1) at (0,0) {$r_1$};
\node[nrn] (E2) at (4,0) {$r_2$};
\node[gab] (G1) at (1.4,1.1) {};
\node[gab] (G2) at (2.6,-1.1) {};
\draw[->,thick] (E1) -- (G1); \draw[inh] (G1) -- (E2);
\draw[->,thick] (E2) -- (G2); \draw[inh] (G2) -- (E1);
\draw[->,thick,blue!60!black] (-1.6,0) node[left,black] {$I_1$} -- (E1);
\draw[->,thick,blue!60!black] (5.6,0) node[right,black] {$I_2$} -- (E2);
\end{tikzpicture}
\caption{二选一。两群\nt{GLUT}神经元（蓝）经\nt{GABA}中介（红）相互抑制。}
\label{fig:wta}
\end{figure}

只要两个神经元都在工作，系统~\eqref{eq:wta} 就是线性的，其行为由两个特征值决定：两个频率同向偏离时为 $-(1+\beta)/\tau$，反向偏离时为 $-(1-\beta)/\tau$。抑制弱时，$\beta<1$，两个特征值都是负的：两个神经元都工作，抑制只是压低弱者。抑制强时，$\beta>1$，第二个特征值变为正：$r_1$ 与 $r_2$ 之间的任何差别都会增长，直到一个神经元完全沉默。

\begin{keyblock}
\begin{proposition}[赢者通吃]
\label{prop:wta}
如果相互抑制强于1，即 $\beta>1$，两群都工作的状态就是不稳定的。网络进入一群工作、另一群沉默的状态。频率为 $r_1=I_1$ 的赢家只要 $I_2<\beta I_1$ 就能保住胜利。
\end{proposition}
\end{keyblock}

我们在模型上验证了这一点。$\beta=0.5$、输入为 $1.0$ 和 $0.9$ 时，两群都工作，频率分别为 $0.73$ 和 $0.53$。$\beta=2$、输入相同时，第二群完全沉默。这种选择带有记忆：如果之后把输入对调，原来的赢家仍然是赢家，因为 $1.0<2\cdot0.9$。

一百群神经元也是一样：每群抑制其余各群，只有一群存活。

\section{存储复位}

第\ref{sec:glutcomputer}章中的存储只能靠疲劳关闭。现在给它加一个\nt{GABA}神经元，由“复位”信号兴奋，并抑制这群神经元。抑制使图~\ref{fig:bistable} 中的曲线 $f(wr+I)$ 下移，上面的交点消失，神经元群熄灭，复位信号结束后仍停留在下方状态。现在存储由一个信号写入、由另一个信号擦除，就像带置位和复位输入的触发器。

更漂亮的做法是：把两个这样的神经元群像图~\ref{fig:wta} 那样用相互抑制连接起来，并给每群加上自反馈。送到其中一群输入端的信号使存储单元进入相应状态，单元记住最后一次决定。这正是由两个交叉耦合的或非门组成的锁存器的直接对应物。

\section{稳定性与节律}

剩下\nt{GLUT}网络的第三种毛病：雪崩。取一群自反馈为 $w_{EE}$ 的\nt{GLUT}神经元，以及一群\nt{GABA}神经元：前者以强度 $w_{IE}$ 兴奋后者，后者以强度 $w_{EI}$ 抑制前者。对频率 $E$ 和 $I$ 有
\begin{equation}
\tau_E\frac{\dd E}{\dd t} = -E + w_{EE}E - w_{EI}I + h, \qquad
\tau_I\frac{\dd I}{\dd t} = -I + w_{IE}E ,
\label{eq:ei}
\end{equation}
其中 $h$ 是外部输入~\cite{WilsonCowan1972}。没有抑制时，若 $w_{EE}>1$，活动会无限增长：每个脉冲引发不止一个后续脉冲。有了抑制，稳态频率
\begin{equation}
E^* \;=\; \frac{h}{1-w_{EE}+w_{EI}w_{IE}}
\label{eq:eistar}
\end{equation}
只要分母为正就是有限的。但这还不够：平衡点还必须是吸引的。对~\eqref{eq:ei} 做线性分析，得到两个条件。

\begin{keyblock}
\begin{proposition}[稳定性条件]
\label{prop:ei}
平衡点~\eqref{eq:eistar} 稳定的条件是抑制
\begin{enumerate}
\item[(i)] \emph{足够强}：$w_{EI}\,w_{IE} > w_{EE}-1$；
\item[(ii)] \emph{足够快}：$\tau_I < \tau_E/(w_{EE}-1)$。
\end{enumerate}
若违反 (i)，活动雪崩式增长。若满足 (i) 而违反 (ii)，抑制来迟，活动开始振荡。
\end{proposition}
\end{keyblock}

条件~(i) 来自系统~\eqref{eq:ei} 矩阵的行列式，条件~(ii) 来自它的迹。第二个条件的含义，任何调过控制器的人都明白：迟到的反馈不能消除偏差，反而会把它激起来。

\begin{figure}[t]
\centering
\begin{tikzpicture}[>=Stealth]
\draw[->,gray!70!black] (0,0) -- (10.9,0) node[below left,black] {\small $t$，ms};
\draw[->,gray!70!black] (0,0) -- (0,3.3) node[left,black] {\small $E$};
\foreach \t in {0,50,100,150}{\draw[gray!70!black] (\t*0.07,0.06) -- (\t*0.07,-0.06) node[below,font=\scriptsize] {\t};}
\draw[very thick,gray!60!black,densely dashed] plot coordinates {(0.000,0.000) (0.035,0.071) (0.070,0.144) (0.105,0.218) (0.140,0.294) (0.175,0.373) (0.210,0.453) (0.245,0.535) (0.280,0.620) (0.315,0.706) (0.350,0.795) (0.385,0.886) (0.420,0.979) (0.455,1.075) (0.490,1.173) (0.525,1.274) (0.560,1.377) (0.595,1.482) (0.630,1.591) (0.665,1.702) (0.700,1.816) (0.735,1.933) (0.770,2.052) (0.805,2.175) (0.840,2.301) (0.875,2.430) (0.910,2.563) (0.945,2.698) (0.980,2.838) (1.015,2.980) (1.050,3.080) (1.085,3.080) (1.120,3.080) (1.155,3.080) (1.190,3.080) (1.225,3.080) (1.260,3.080) (1.295,3.080) (1.330,3.080) (1.365,3.080) (1.400,3.080) (1.435,3.080) (1.470,3.080) (1.505,3.080) (1.540,3.080) (1.575,3.080) (1.610,3.080) (1.645,3.080) (1.680,3.080) (1.715,3.080) (1.750,3.080) (1.785,3.080) (1.820,3.080) (1.855,3.080) (1.890,3.080) (1.925,3.080) (1.960,3.080) (1.995,3.080) (2.030,3.080) (2.065,3.080) (2.100,3.080) (2.135,3.080) (2.170,3.080) (2.205,3.080) (2.240,3.080) (2.275,3.080) (2.310,3.080) (2.345,3.080) (2.380,3.080) (2.415,3.080) (2.450,3.080) (2.485,3.080) (2.520,3.080) (2.555,3.080) (2.590,3.080) (2.625,3.080) (2.660,3.080) (2.695,3.080) (2.730,3.080) (2.765,3.080) (2.800,3.080) (2.835,3.080) (2.870,3.080) (2.905,3.080) (2.940,3.080) (2.975,3.080) (3.010,3.080) (3.045,3.080) (3.080,3.080) (3.115,3.080) (3.150,3.080) (3.185,3.080) (3.220,3.080) (3.255,3.080) (3.290,3.080) (3.325,3.080) (3.360,3.080) (3.395,3.080) (3.430,3.080) (3.465,3.080) (3.500,3.080) (3.535,3.080) (3.570,3.080) (3.605,3.080) (3.640,3.080) (3.675,3.080) (3.710,3.080) (3.745,3.080) (3.780,3.080) (3.815,3.080) (3.850,3.080) (3.885,3.080) (3.920,3.080) (3.955,3.080) (3.990,3.080) (4.025,3.080) (4.060,3.080) (4.095,3.080) (4.130,3.080) (4.165,3.080) (4.200,3.080) (4.235,3.080) (4.270,3.080) (4.305,3.080) (4.340,3.080) (4.375,3.080) (4.410,3.080) (4.445,3.080) (4.480,3.080) (4.515,3.080) (4.550,3.080) (4.585,3.080) (4.620,3.080) (4.655,3.080) (4.690,3.080) (4.725,3.080) (4.760,3.080) (4.795,3.080) (4.830,3.080) (4.865,3.080) (4.900,3.080) (4.935,3.080) (4.970,3.080) (5.005,3.080) (5.040,3.080) (5.075,3.080) (5.110,3.080) (5.145,3.080) (5.180,3.080) (5.215,3.080) (5.250,3.080) (5.285,3.080) (5.320,3.080) (5.355,3.080) (5.390,3.080) (5.425,3.080) (5.460,3.080) (5.495,3.080) (5.530,3.080) (5.565,3.080) (5.600,3.080) (5.635,3.080) (5.670,3.080) (5.705,3.080) (5.740,3.080) (5.775,3.080) (5.810,3.080) (5.845,3.080) (5.880,3.080) (5.915,3.080) (5.950,3.080) (5.985,3.080) (6.020,3.080) (6.055,3.080) (6.090,3.080) (6.125,3.080) (6.160,3.080) (6.195,3.080) (6.230,3.080) (6.265,3.080) (6.300,3.080) (6.335,3.080) (6.370,3.080) (6.405,3.080) (6.440,3.080) (6.475,3.080) (6.510,3.080) (6.545,3.080) (6.580,3.080) (6.615,3.080) (6.650,3.080) (6.685,3.080) (6.720,3.080) (6.755,3.080) (6.790,3.080) (6.825,3.080) (6.860,3.080) (6.895,3.080) (6.930,3.080) (6.965,3.080) (7.000,3.080) (7.035,3.080) (7.070,3.080) (7.105,3.080) (7.140,3.080) (7.175,3.080) (7.210,3.080) (7.245,3.080) (7.280,3.080) (7.315,3.080) (7.350,3.080) (7.385,3.080) (7.420,3.080) (7.455,3.080) (7.490,3.080) (7.525,3.080) (7.560,3.080) (7.595,3.080) (7.630,3.080) (7.665,3.080) (7.700,3.080) (7.735,3.080) (7.770,3.080) (7.805,3.080) (7.840,3.080) (7.875,3.080) (7.910,3.080) (7.945,3.080) (7.980,3.080) (8.015,3.080) (8.050,3.080) (8.085,3.080) (8.120,3.080) (8.155,3.080) (8.190,3.080) (8.225,3.080) (8.260,3.080) (8.295,3.080) (8.330,3.080) (8.365,3.080) (8.400,3.080) (8.435,3.080) (8.470,3.080) (8.505,3.080) (8.540,3.080) (8.575,3.080) (8.610,3.080) (8.645,3.080) (8.680,3.080) (8.715,3.080) (8.750,3.080) (8.785,3.080) (8.820,3.080) (8.855,3.080) (8.890,3.080) (8.925,3.080) (8.960,3.080) (8.995,3.080) (9.030,3.080) (9.065,3.080) (9.100,3.080) (9.135,3.080) (9.170,3.080) (9.205,3.080) (9.240,3.080) (9.275,3.080) (9.310,3.080) (9.345,3.080) (9.380,3.080) (9.415,3.080) (9.450,3.080) (9.485,3.080) (9.520,3.080) (9.555,3.080) (9.590,3.080) (9.625,3.080) (9.660,3.080) (9.695,3.080) (9.730,3.080) (9.765,3.080) (9.800,3.080) (9.835,3.080) (9.870,3.080) (9.905,3.080) (9.940,3.080) (9.975,3.080) (10.010,3.080) (10.045,3.080) (10.080,3.080) (10.115,3.080) (10.150,3.080) (10.185,3.080) (10.220,3.080) (10.255,3.080) (10.290,3.080) (10.325,3.080) (10.360,3.080) (10.395,3.080) (10.430,3.080) (10.465,3.080) (10.500,3.080)};
\draw[very thick,blue!60!black] plot coordinates {(0.000,0.000) (0.035,0.071) (0.070,0.141) (0.105,0.211) (0.140,0.279) (0.175,0.344) (0.210,0.406) (0.245,0.465) (0.280,0.519) (0.315,0.570) (0.350,0.616) (0.385,0.658) (0.420,0.697) (0.455,0.731) (0.490,0.763) (0.525,0.790) (0.560,0.815) (0.595,0.836) (0.630,0.855) (0.665,0.871) (0.700,0.886) (0.735,0.898) (0.770,0.908) (0.805,0.916) (0.840,0.924) (0.875,0.929) (0.910,0.934) (0.945,0.938) (0.980,0.941) (1.015,0.943) (1.050,0.945) (1.085,0.946) (1.120,0.946) (1.155,0.947) (1.190,0.947) (1.225,0.947) (1.260,0.946) (1.295,0.946) (1.330,0.945) (1.365,0.944) (1.400,0.944) (1.435,0.943) (1.470,0.942) (1.505,0.941) (1.540,0.941) (1.575,0.940) (1.610,0.939) (1.645,0.939) (1.680,0.938) (1.715,0.937) (1.750,0.937) (1.785,0.936) (1.820,0.936) (1.855,0.936) (1.890,0.935) (1.925,0.935) (1.960,0.935) (1.995,0.934) (2.030,0.934) (2.065,0.934) (2.100,0.934) (2.135,0.934) (2.170,0.934) (2.205,0.934) (2.240,0.933) (2.275,0.933) (2.310,0.933) (2.345,0.933) (2.380,0.933) (2.415,0.933) (2.450,0.933) (2.485,0.933) (2.520,0.933) (2.555,0.933) (2.590,0.933) (2.625,0.933) (2.660,0.933) (2.695,0.933) (2.730,0.933) (2.765,0.933) (2.800,0.933) (2.835,0.933) (2.870,0.933) (2.905,0.933) (2.940,0.933) (2.975,0.933) (3.010,0.933) (3.045,0.933) (3.080,0.933) (3.115,0.933) (3.150,0.933) (3.185,0.933) (3.220,0.933) (3.255,0.933) (3.290,0.933) (3.325,0.933) (3.360,0.933) (3.395,0.933) (3.430,0.933) (3.465,0.933) (3.500,0.933) (3.535,0.933) (3.570,0.933) (3.605,0.933) (3.640,0.933) (3.675,0.933) (3.710,0.933) (3.745,0.933) (3.780,0.933) (3.815,0.933) (3.850,0.933) (3.885,0.933) (3.920,0.933) (3.955,0.933) (3.990,0.933) (4.025,0.933) (4.060,0.933) (4.095,0.933) (4.130,0.933) (4.165,0.933) (4.200,0.933) (4.235,0.933) (4.270,0.933) (4.305,0.933) (4.340,0.933) (4.375,0.933) (4.410,0.933) (4.445,0.933) (4.480,0.933) (4.515,0.933) (4.550,0.933) (4.585,0.933) (4.620,0.933) (4.655,0.933) (4.690,0.933) (4.725,0.933) (4.760,0.933) (4.795,0.933) (4.830,0.933) (4.865,0.933) (4.900,0.933) (4.935,0.933) (4.970,0.933) (5.005,0.933) (5.040,0.933) (5.075,0.933) (5.110,0.933) (5.145,0.933) (5.180,0.933) (5.215,0.933) (5.250,0.933) (5.285,0.933) (5.320,0.933) (5.355,0.933) (5.390,0.933) (5.425,0.933) (5.460,0.933) (5.495,0.933) (5.530,0.933) (5.565,0.933) (5.600,0.933) (5.635,0.933) (5.670,0.933) (5.705,0.933) (5.740,0.933) (5.775,0.933) (5.810,0.933) (5.845,0.933) (5.880,0.933) (5.915,0.933) (5.950,0.933) (5.985,0.933) (6.020,0.933) (6.055,0.933) (6.090,0.933) (6.125,0.933) (6.160,0.933) (6.195,0.933) (6.230,0.933) (6.265,0.933) (6.300,0.933) (6.335,0.933) (6.370,0.933) (6.405,0.933) (6.440,0.933) (6.475,0.933) (6.510,0.933) (6.545,0.933) (6.580,0.933) (6.615,0.933) (6.650,0.933) (6.685,0.933) (6.720,0.933) (6.755,0.933) (6.790,0.933) (6.825,0.933) (6.860,0.933) (6.895,0.933) (6.930,0.933) (6.965,0.933) (7.000,0.933) (7.035,0.933) (7.070,0.933) (7.105,0.933) (7.140,0.933) (7.175,0.933) (7.210,0.933) (7.245,0.933) (7.280,0.933) (7.315,0.933) (7.350,0.933) (7.385,0.933) (7.420,0.933) (7.455,0.933) (7.490,0.933) (7.525,0.933) (7.560,0.933) (7.595,0.933) (7.630,0.933) (7.665,0.933) (7.700,0.933) (7.735,0.933) (7.770,0.933) (7.805,0.933) (7.840,0.933) (7.875,0.933) (7.910,0.933) (7.945,0.933) (7.980,0.933) (8.015,0.933) (8.050,0.933) (8.085,0.933) (8.120,0.933) (8.155,0.933) (8.190,0.933) (8.225,0.933) (8.260,0.933) (8.295,0.933) (8.330,0.933) (8.365,0.933) (8.400,0.933) (8.435,0.933) (8.470,0.933) (8.505,0.933) (8.540,0.933) (8.575,0.933) (8.610,0.933) (8.645,0.933) (8.680,0.933) (8.715,0.933) (8.750,0.933) (8.785,0.933) (8.820,0.933) (8.855,0.933) (8.890,0.933) (8.925,0.933) (8.960,0.933) (8.995,0.933) (9.030,0.933) (9.065,0.933) (9.100,0.933) (9.135,0.933) (9.170,0.933) (9.205,0.933) (9.240,0.933) (9.275,0.933) (9.310,0.933) (9.345,0.933) (9.380,0.933) (9.415,0.933) (9.450,0.933) (9.485,0.933) (9.520,0.933) (9.555,0.933) (9.590,0.933) (9.625,0.933) (9.660,0.933) (9.695,0.933) (9.730,0.933) (9.765,0.933) (9.800,0.933) (9.835,0.933) (9.870,0.933) (9.905,0.933) (9.940,0.933) (9.975,0.933) (10.010,0.933) (10.045,0.933) (10.080,0.933) (10.115,0.933) (10.150,0.933) (10.185,0.933) (10.220,0.933) (10.255,0.933) (10.290,0.933) (10.325,0.933) (10.360,0.933) (10.395,0.933) (10.430,0.933) (10.465,0.933) (10.500,0.933)};
\draw[very thick,red!70!black] plot coordinates {(0.000,0.000) (0.035,0.071) (0.070,0.143) (0.105,0.217) (0.140,0.293) (0.175,0.370) (0.210,0.449) (0.245,0.528) (0.280,0.609) (0.315,0.691) (0.350,0.774) (0.385,0.858) (0.420,0.943) (0.455,1.028) (0.490,1.114) (0.525,1.200) (0.560,1.287) (0.595,1.374) (0.630,1.462) (0.665,1.549) (0.700,1.636) (0.735,1.724) (0.770,1.810) (0.805,1.897) (0.840,1.983) (0.875,2.068) (0.910,2.153) (0.945,2.237) (0.980,2.320) (1.015,2.402) (1.050,2.483) (1.085,2.563) (1.120,2.641) (1.155,2.717) (1.190,2.792) (1.225,2.866) (1.260,2.937) (1.295,3.007) (1.330,3.075) (1.365,3.080) (1.400,3.080) (1.435,3.080) (1.470,3.080) (1.505,3.080) (1.540,3.080) (1.575,3.080) (1.610,3.080) (1.645,3.080) (1.680,3.080) (1.715,3.080) (1.750,3.080) (1.785,3.080) (1.820,3.080) (1.855,3.080) (1.890,3.080) (1.925,3.080) (1.960,3.080) (1.995,3.080) (2.030,3.080) (2.065,3.080) (2.100,3.080) (2.135,3.080) (2.170,3.080) (2.205,3.080) (2.240,3.080) (2.275,3.080) (2.310,3.080) (2.345,3.080) (2.380,3.080) (2.415,3.080) (2.450,3.080) (2.485,3.080) (2.520,3.080) (2.555,3.080) (2.590,3.080) (2.625,3.080) (2.660,3.080) (2.695,3.080) (2.730,3.080) (2.765,3.080) (2.800,3.080) (2.835,3.052) (2.870,2.973) (2.905,2.892) (2.940,2.807) (2.975,2.719) (3.010,2.629) (3.045,2.535) (3.080,2.438) (3.115,2.339) (3.150,2.238) (3.185,2.133) (3.220,2.029) (3.255,1.930) (3.290,1.836) (3.325,1.747) (3.360,1.661) (3.395,1.580) (3.430,1.503) (3.465,1.430) (3.500,1.360) (3.535,1.294) (3.570,1.231) (3.605,1.171) (3.640,1.113) (3.675,1.059) (3.710,1.007) (3.745,0.958) (3.780,0.911) (3.815,0.867) (3.850,0.825) (3.885,0.784) (3.920,0.746) (3.955,0.710) (3.990,0.675) (4.025,0.642) (4.060,0.611) (4.095,0.581) (4.130,0.553) (4.165,0.526) (4.200,0.500) (4.235,0.476) (4.270,0.452) (4.305,0.430) (4.340,0.409) (4.375,0.389) (4.410,0.370) (4.445,0.352) (4.480,0.335) (4.515,0.319) (4.550,0.303) (4.585,0.288) (4.620,0.274) (4.655,0.261) (4.690,0.248) (4.725,0.236) (4.760,0.225) (4.795,0.214) (4.830,0.203) (4.865,0.193) (4.900,0.184) (4.935,0.175) (4.970,0.166) (5.005,0.158) (5.040,0.151) (5.075,0.143) (5.110,0.136) (5.145,0.130) (5.180,0.123) (5.215,0.117) (5.250,0.111) (5.285,0.106) (5.320,0.101) (5.355,0.096) (5.390,0.091) (5.425,0.087) (5.460,0.083) (5.495,0.079) (5.530,0.075) (5.565,0.071) (5.600,0.068) (5.635,0.064) (5.670,0.061) (5.705,0.058) (5.740,0.055) (5.775,0.053) (5.810,0.050) (5.845,0.048) (5.880,0.045) (5.915,0.043) (5.950,0.041) (5.985,0.039) (6.020,0.037) (6.055,0.035) (6.090,0.034) (6.125,0.032) (6.160,0.030) (6.195,0.029) (6.230,0.027) (6.265,0.026) (6.300,0.025) (6.335,0.024) (6.370,0.022) (6.405,0.021) (6.440,0.021) (6.475,0.022) (6.510,0.024) (6.545,0.027) (6.580,0.032) (6.615,0.037) (6.650,0.044) (6.685,0.052) (6.720,0.061) (6.755,0.071) (6.790,0.083) (6.825,0.095) (6.860,0.109) (6.895,0.124) (6.930,0.140) (6.965,0.157) (7.000,0.176) (7.035,0.195) (7.070,0.215) (7.105,0.237) (7.140,0.260) (7.175,0.283) (7.210,0.308) (7.245,0.333) (7.280,0.360) (7.315,0.388) (7.350,0.416) (7.385,0.445) (7.420,0.476) (7.455,0.507) (7.490,0.539) (7.525,0.571) (7.560,0.604) (7.595,0.638) (7.630,0.673) (7.665,0.708) (7.700,0.744) (7.735,0.781) (7.770,0.818) (7.805,0.855) (7.840,0.893) (7.875,0.931) (7.910,0.969) (7.945,1.008) (7.980,1.047) (8.015,1.086) (8.050,1.126) (8.085,1.165) (8.120,1.204) (8.155,1.244) (8.190,1.283) (8.225,1.322) (8.260,1.362) (8.295,1.400) (8.330,1.439) (8.365,1.477) (8.400,1.515) (8.435,1.553) (8.470,1.590) (8.505,1.626) (8.540,1.662) (8.575,1.698) (8.610,1.733) (8.645,1.767) (8.680,1.800) (8.715,1.832) (8.750,1.864) (8.785,1.895) (8.820,1.924) (8.855,1.953) (8.890,1.981) (8.925,2.007) (8.960,2.033) (8.995,2.057) (9.030,2.080) (9.065,2.102) (9.100,2.123) (9.135,2.142) (9.170,2.160) (9.205,2.177) (9.240,2.192) (9.275,2.205) (9.310,2.217) (9.345,2.228) (9.380,2.237) (9.415,2.245) (9.450,2.251) (9.485,2.255) (9.520,2.258) (9.555,2.259) (9.590,2.259) (9.625,2.256) (9.660,2.253) (9.695,2.247) (9.730,2.240) (9.765,2.231) (9.800,2.220) (9.835,2.208) (9.870,2.194) (9.905,2.178) (9.940,2.160) (9.975,2.141) (10.010,2.120) (10.045,2.098) (10.080,2.074) (10.115,2.048) (10.150,2.021) (10.185,1.992) (10.220,1.961) (10.255,1.929) (10.290,1.895) (10.325,1.860) (10.360,1.824) (10.395,1.786) (10.430,1.746) (10.465,1.706) (10.500,1.663)};

\node[font=\small,gray!40!black,anchor=west] at (0.9,3.15) {无抑制：雪崩};
\node[font=\small,blue!60!black,anchor=west] at (7.4,1.25) {快速抑制};
\node[font=\small,red!70!black,anchor=west] at (7.4,2.55) {慢速抑制};
\end{tikzpicture}
\caption{接通输入后，反馈为 $w_{EE}=1.5$、$\tau_E=10\ms$ 的\nt{GLUT}神经元群的频率。无抑制时（虚线）活动无限增长。有强度 $w_{EI}w_{IE}=2$、$\tau_I=2\ms$ 的抑制时（蓝色曲线），活动达到 $E^*=h/(1-1.5+2)=0.67h$ 的水平。抑制强度相同但较慢，$\tau_I=30\ms$ 时（红色曲线），活动发生振荡。}
\label{fig:ei}
\end{figure}

一般认为，皮层正是工作在这种状态：它自身的反馈强到足以在没有抑制时进入雪崩，只靠快速抑制把它维持在工作点上~\cite{Tsodyks1997isn}。

这种状态有一个反直觉的特征，可以从外部辨认出来~\cite{Tsodyks1997isn}。在~\eqref{eq:ei} 中给\nt{GABA}神经元群直接加上外部输入 $h_I$。稳态频率变为
\begin{equation}
E^* \;=\; \frac{h-w_{EI}\,h_I}{D},\qquad
I^* \;=\; \frac{w_{IE}\,h+(1-w_{EE})\,h_I}{D},\qquad
D \;=\; 1-w_{EE}+w_{EI}w_{IE} .
\label{eq:isn}
\end{equation}
若 $w_{EE}>1$，第二个公式中 $h_I$ 的系数为负：推一下\nt{GABA}神经元，它们的频率反而\emph{下降}。原因在于环路。额外的抑制压低了\nt{GLUT}神经元群，后者对\nt{GABA}神经元的兴奋减少，于是\nt{GABA}神经元失去的比得到的更多。按图~\ref{fig:ei} 的数值（$w_{EE}=1.5$，$w_{EI}w_{IE}=2$，$D=1.5$），每增加一个单位的输入，$I^*$ 就降低三分之一个单位。皮层中确实发现了这个特征：当视皮层神经元的响应被周围刺激抑制时，它们收到的抑制性电流不增反降，与兴奋性电流一同下降~\cite{Ozeki2009}。后来在皮层的不同区域和不同层中都发现了同样的特征~\cite{Sanzeni2020}。

违反条件~(ii) 时的振荡在大脑中同样存在。“\nt{GLUT}兴奋\nt{GABA}，\nt{GABA}抑制\nt{GLUT}”这个带几毫秒延迟的环路，产生周期约 $12$--$50\ms$（频率 $20$--$80$\,Hz）的节律，皮层中这样的快速节律众所周知~\cite{Cardin2009,Buzsaki2012}。这不是计算机的时钟：节律是局部的，只在网络活跃时才出现。但在几个周期之内，它把时间划分成神经元可以发放的窗口。

\section{小结}

\begin{table}[t]
\centering
\small
\begin{tabular}{>{\raggedright}p{3.3cm}>{\raggedright}p{4.4cm}>{\raggedright\arraybackslash}p{4.4cm}}
\toprule
& 只有\nt{GLUT} & \nt{GLUT} + \nt{GABA} \\
\midrule
非 & 只能借助“接通—断开”传感器 & 否决 $a\wedge\bar b$；有背景活动时可实现非 \\
每比特导线数 & 两根 & 一根 \\
运动方向 & 无 & 带延迟的否决 \\
增益调节 & 无 & 除法：分流加背景噪声 \\
符合窗口 & $\sim\tau$ & 被抑制收窄 \\
多选一 & 无 & 相互抑制，$\beta>1$ \\
存储复位 & 只能靠疲劳 & 经\nt{GABA}用信号复位 \\
稳定性 & 只能靠疲劳 & 快速负反馈 \\
节律 & 不需要 & 抑制慢时出现 \\
\bottomrule
\end{tabular}
\caption{\nt{GABA}给\nt{GLUT}神经元网络带来了什么。}
\label{tab:gaba}
\end{table}

\nt{GABA}几乎没有给网络增加新的\emph{计算}——这些用双线传感器也能算出来——它增加的是\emph{控制}（表~\ref{tab:gaba}）：选择、复位、增益调节和稳定性。大脑中的兴奋承载数据，而抑制决定哪些数据通过、何时通过、以多大音量通过。对皮层中每四五个神经元中的一个而言，这是非常合理的分工。

\section{习题}

\begin{exercise}[\lvlA]
\label{ex:03:1}
\emph{分流的数值。}$E_E=70\mV$。由~\eqref{eq:shunt} 求 $g_E=g_L$ 和 $4g_L$ 时的 $V_\infty$，分别在无抑制和有分流 $g_I=2g_L$ 的情况下。一个在 $g_E=g_L$ 时与分流扣除等量电压的减法，在 $g_E=4g_L$ 时会给出多大的 $V_\infty$？分流使符合窗口收窄几倍？
\end{exercise}

\begin{exercise}[\lvlB]
\label{ex:03:2}
\emph{推导稳定性条件。}求线性系统~\eqref{eq:ei} 矩阵的迹和行列式，并由此推出命题~\ref{prop:ei} 的两个条件。在条件~(ii) 的边界上，平衡点以振荡方式失稳。按图~\ref{fig:ei} 的数值求这一振荡的周期。
\end{exercise}

\begin{exercise}[\lvlB]
\label{ex:03:3}
\emph{由延迟产生节律。}把模型~\eqref{eq:ei} 中的慢抑制换成带延迟 $d$ 的快抑制：$\tau_E\dot E=-E(t)-GE(t-d)+h$。$\tau_E=10\ms$ 时，对 $G=2$ 和 $G=5$ 求产生振荡的最小延迟 $d_c$ 及振荡周期。与习题~\ref{ex:03:2} 不同，它们是否落在皮层快速节律的 $12$--$50\ms$ 范围内？
\end{exercise}

\begin{exercise}[\lvlB]
\label{ex:03:4}
\emph{建模：带记忆的选择。}在 $\beta=2$、$\tau=1$ 时对~\eqref{eq:wta} 积分。(a)~$I_1=1$，把 $I_2$ 从 $0$ 缓慢升到 $3$ 再降回来：网络在哪些 $I_2$ 处切换？(b)~当输入差 $\varepsilon$ 缩小到十分之一时，从静默出发的决策时间增加多少？
\end{exercise}

\begin{exercise}[\lvlB]
\label{ex:03:5}
\emph{设计：速度带检测器。}图~\ref{fig:direction} 的检测器在从 $A$ 到 $B$ 的运动、渡越时间为 $5$--$20\ms$ 时必须沉默。延迟为 $d_k$ 的\nt{GABA}中介在 $A$ 发放后的区间 $[d_k,\,d_k+5\ms]$ 内否决 $Y$；来自 $B$ 的兴奋立即到达。需要几个中介、延迟各为多少？
\end{exercise}

\begin{exercise}[\lvlB]
\label{ex:03:6}
\emph{否决与背景。}在“每个输入一个\nt{GABA}中介、每个 $f=1$ 的输入组合一个否决、再加一个总的或”的通用方案中，奇偶校验 $x\oplus y\oplus z$ 需要多少个神经元？用两个直接接收输入的神经元（一个\nt{GABA}、一个\nt{GLUT}）实现它，给出权重和阈值。
\end{exercise}

\begin{exercise}[\lvlC]
\label{ex:03:7}
\emph{设计：百里挑一。}$100$ 群\nt{GLUT}中的每一群都必须同等地抑制所有其他群，但不抑制自己。线性\nt{GABA}中介由神经元群兴奋并抑制它们；神经元群可以相互兴奋，但不能自我兴奋。中介最少要几个？如果完全没有直接的兴奋性连接呢？
\end{exercise}

\begin{exercise}[\lvlC]
\label{ex:03:8}
\emph{研究：抑制何时反常。}在图~\ref{fig:ei} 的模型中（$w_{EI}=2$，$w_{IE}=1$），\nt{GLUT}群的传递函数改为 $f(u)=[u]_+^2$，\nt{GABA}群仍是线性的。输入超过哪个值 $h_c$ 时，给\nt{GABA}群增加的输入反而降低它自己的频率？用模拟加以验证。
\end{exercise}
